\appendix
\section{Appendix}


\subsection{Composition of power series and Faa di Bruno formula}

In this appendix, we state the necessary definitions and results on composition of power series, referencing \cite[Chapter~11]{Charalambides}.

\begin{Definition}[{\cite{Charalambides}}]
The polynomials $B_n(x_1,x_2,\dots,x_n)$, $B_{n,k}(x_1,x_2,\dots,x_n)$ in the variables~${x_1,x_2,\dots,x_n}$, defined by the sum
\begin{gather}
B_{n}=\sum_{ k_1+2k_2+3k_3+\dots +nk_{n}=n} \frac{n!}{k_1!(1!)k_2!(2!)^{k_2}\cdots k_n!(n!)^{k_n}} x_1^{k_1}x_2^{k_2}\cdots x_n^{k_n},\nonumber\\
B_{n,k}=\sum_{\substack{ k_1+\dots +k_n=k\\ k_1+2k_2+3k_3+\dots +nk_{n}=n}} \frac{n!}{k_1!(1!)k_2!(2!)^{k_2}\cdots k_n!(n!)^{k_n}} x_1^{k_1}x_2^{k_2}\cdots x_n^{k_n}\label{Bell polynomials}
\end{gather}
 are called exponential Bell partition polynomial and partial Bell partition polynomial respectively.
\end{Definition}


\begin{Theorem}[Faa di Bruno formula, \cite{Charalambides}]
Let $f(u)$ and $g(t)$ be two functions of real variables for which all the derivatives
\begin{gather*}
g_r=\left[\frac{{\rm d}^rg(t)}{{\rm d}t^r}\right]_{t=a}, \qquad r=0,1,\dots, \qquad f_k=\left[\frac{{\rm d}^kf(u)}{{\rm d}u^k}\right]_{u=(g(a)}, \qquad k=0,1,\dots,
\end{gather*}
exist. Then the derivatives of the composite function $h(t)= f(g(t))$,
\begin{gather*}
h_n=\left[\frac{{\rm d}^nh(u)}{{\rm d}t^n}\right]_{t=a}, \qquad n=0,1,\dots,
\end{gather*}
are given by
\begin{gather}
h_n=\sum_{k=0}^n f_kB_{n,k}(g_1,g_2,\dots,g_n)=B_n(fg_1,fg_2,\dots,fg_n).\label{Faa di Bruno}
\end{gather}
\end{Theorem}


\begin{Definition}[{\cite{Charalambides}}]
The polynomial $C_{n,s}=C_{n,s}(x_1,x_2,\dots,x_n)$ in the variables $x_1,x_2,\dots,\allowbreak x_n$, defined for a real (or complex) number s by the sum
\begin{gather*}
C_{n,s}=\sum_{ k_1+2k_2+3k_3+\dots +nk_{n}=n} \frac{n!}{k_1!(1!)k_2!(2!)^{k_2}\cdots k_n!(n!)^{k_n}}{{s}\choose{k}} x_1^{k_1}x_2^{k_2}\cdots x_n^{k_n}
\end{gather*}
 is called potential partition polynomial,
 where
 \begin{gather*}
 {{s}\choose{k}}=\frac{s(s-1)(s-2)\cdots (s-k+1))}{k!}.
 \end{gather*}
\end{Definition}


\begin{Theorem}[\cite{Charalambides}]\label{theorem potential partition polynomials}
The generating function of the potential partition polynomials $C_{n,s}(x_1,\dots,\allowbreak x_n)$, $n=0,1,\dots$, for fixed $k$, is given by
 \[%\label{generating function of Cst}
C_{s}(t)=\sum_{n=0}^{\infty} C_{n,s}(x_1,\dots, x_n)\frac{t^n}{n!}=[ 1+(g(t)-x_0) ]^s ,
\]
where
$g(t)=\sum_{r=0}^{\infty} x_r \frac{t^r}{r!}$.
\end{Theorem}

\subsection{Canonical coordinates as functions of Gromov--Witten invariants}\label{Canonical coordinates as functions of Gromov Witten invariants}

\begin{Lemma}\label{lemma of exp 2pi 3}
Let be \smash{$\zeta={\rm e}^{\frac{2\pi {\rm i}}{3}}$} and
\begin{gather*}
b_{k_1,k_2}:=\zeta^{k_1}+\zeta^{k_1+2k_2}+\zeta^{2k_1},\qquad
c_{n_1,n_2,n_3}=\zeta^{n_1+2n_2},\\
\tilde c_{n_1,n_2,n_3}:=c_{n_1,n_2,n_3}+c_{n_1,n_3,n_2}+c_{n_3,n_2,n_1}+c_{n_3,n_1,n_2}+c_{n_2,n_1,n_3}+c_{n_2,n_3,n_1}
\end{gather*}
for $k,k_1,k_2 k_3 \in \mathbb{Z}$. Then the following identities holds:
\begin{gather}
\zeta^{k}+\zeta^{2k}+\zeta^{3k}=1+2\cos\left(\frac{2\pi k}{3} \right),\qquad
b_{k,k}=1+2\cos\left(\frac{2\pi k}{3} \right),\qquad
c_{k,k,k}=1,
\nonumber\\
 b_{k_1-k_2,k_2}+b_{k_2,k_1-k_2}=
\begin{cases}
 6\cos\bigl( \frac{2\pi k_2}{3}\bigr)& \text{if } k_1\in 3\mathbb{Z},\\
 0 & \text{otherwise},
\end{cases}\nonumber
\\
\tilde c_{k_1-k_2-k_3,k_2,k_3}=
\begin{cases}
 6\cos\bigl( \frac{2\pi (k_2+2k_3)}{3}\bigr)& \text{if } k_1\in 3\mathbb{Z},\\
 0 & \text{otherwise}.
\end{cases}\label{exp 2pi 3 id}
\end{gather}
\end{Lemma}

\begin{proof}
\[
\zeta^{k}+\zeta^{2k}+\zeta^{3k}=\zeta^{k}+\zeta^{2k}+1=2\left(\frac{{\rm e}^{\frac{2\pi {\rm i}k}{3}}+{\rm e}^{\frac{-2\pi {\rm i}k}{3}}}{2} \right)+1=1+2\cos\left(\frac{2\pi k}{3} \right).
\]

Define $b_{k_1,k_2}$,
$
b_{k_1,k_2}:=\zeta^{k_1}+\zeta^{k_1+2k_2}+\zeta^{2k_1}$.
Then,
\[
b_{k,k}=\zeta^{k}+\zeta^{3k}+\zeta^{2k}=1+2\cos\left(\frac{2\pi k}{3} \right).
\]
Moreover,
\[
b_{k_1-k_2,k_2}=\zeta^{k_1-k_2}+\zeta^{k_1+k_2}+\zeta^{2(k_1-k_2)},\qquad
b_{k_2,k_1-k_2}=\zeta^{k_2}+\zeta^{2k_1-k_2}+\zeta^{2k_2}.
\]

If $k_1 \in 3\mathbb{Z}$,
\begin{align*}
b_{k_1-k_2,k_2}+b_{k_2,k_1-k_2}&=\zeta^{-k_2}+\zeta^{k_2}+\zeta^{-2k_2}+\zeta^{k_2}+\zeta^{-k_2}+\zeta^{2k_2}\\
&=\zeta^{2k_2}+\zeta^{k_2}+\zeta^{k_2}+\zeta^{k_2}+\zeta^{2k_2}+\zeta^{2k_2}\\
&=3\bigl( \zeta^{2k_2}+\zeta^{k_2}\bigr)=6\cos\left( \frac{2\pi k_2}{3}\right).
\end{align*}

If $k_1 \in 1+3\mathbb{Z}$,
\begin{align*}
b_{k_1-k_2,k_2}+b_{k_2,k_1-k_2}&=\zeta^{1-k_2}+\zeta^{k_2+1}+\zeta^{2-2k_2}+\zeta^{k_2}+\zeta^{2-k_2}+\zeta^{2k_2}\\
&=\zeta^{1+2k_2}+\zeta^{k_2+1}+\zeta^{k_2+2}+\zeta^{k_2}+\zeta^{2(k_2+1)}+\zeta^{2k_2}\\
&=2\cos\left( \frac{2\pi (1-k_2)}{3}\right)+2\cos\left( \frac{2\pi k_2}{3}\right)+2\cos\left( \frac{2\pi (1+k_2)}{3}\right)\\
&=0.
\end{align*}

If $k_1 \in 2+3\mathbb{Z}$,
\begin{align*}
b_{k_1-k_2,k_2}+b_{k_2,k_1-k_2}&=\zeta^{2-k_2}+\zeta^{k_2+2}+\zeta^{1-2k_2}+\zeta^{k_2}+\zeta^{1-k_2}+\zeta^{2k_2}\\
&=\zeta^{2-2k_2}+\zeta^{k_2+2}+\zeta^{2(2-k_2)}+\zeta^{k_2}+\zeta^{2(k_2+2)}+\zeta^{2k_2}\\
&=2\cos\left( \frac{2\pi (2-k_2)}{3}\right)+2\cos\left( \frac{2\pi k_2}{3}\right)+2\cos\left( \frac{2\pi (2+k_2)}{3}\right)\\
&=0.
\end{align*}

Let be the tensor defined by
$c_{n_1,n_2,n_3}=\zeta^{n_1+2n_2}$.
It is straightforward the following
$c_{n,n,n}=\zeta^{3n}=1$.

Moreover, setting
\begin{gather*}
\tilde c_{n_1,n_2,n_3}:=c_{n_1,n_2,n_3}+c_{n_1,n_3,n_2}+c_{n_3,n_2,n_1}+c_{n_3,n_1,n_2}+c_{n_2,n_1,n_3}+c_{n_2,n_3,n_1},
\end{gather*}
we have
\begin{gather*}
\tilde c_{n_1,n_2,n_3}:=\zeta^{n_1+2n_2}+\zeta^{n_1+2n_3}+\zeta^{n_3+2n_2}+\zeta^{n_3+2n_1}+\zeta^{n_2+2n_1}+\zeta^{n_2+2n_3}.
\end{gather*}
In particular,
\begin{align*}
\tilde c_{n_1-n_2-n_3,n_2,n_3}={}&\zeta^{n_1+n_2-n_3}+\zeta^{n_1-n_2+n_3}+\zeta^{n_3+2n_2}\\
&+\zeta^{-n_3+2n_1-2n_2}+\zeta^{-n_2+2n_1-2n_3}+\zeta^{n_2+2n_3}.
\end{align*}

If $n_1 \in 3\mathbb{Z}$,
\begin{align*}
\tilde c_{n_1-n_2-n_3,n_2,n_3}&=\zeta^{n_2-n_3}+\zeta^{-n_2+n_3}+\zeta^{n_3+2n_2}+\zeta^{-n_3-2n_2}+\zeta^{-n_2-2n_3}+\zeta^{n_2+2n_3}\\
&=\zeta^{n_2+2n_3}+\zeta^{2n_2+n_3}+\zeta^{n_3+2n_2}+\zeta^{2n_3+n_2}+\zeta^{2n_2+n_3}+\zeta^{n_2+2n_3}\\
&=3\bigl(\zeta^{n_2+2n_3}+\zeta^{2n_2+n_3}\bigr)
=3\bigl(\zeta^{n_2+2n_3}+\zeta^{-n_2-2n_3}\bigr)\\
&=6\cos\left( \frac{2\pi \left(n_2+2n_3 \right)}{3} \right).
\end{align*}

If $n_1 \in 1+3\mathbb{Z}$,
\begin{align*}
\tilde c_{n_1-n_2-n_3,n_2,n_3}={}&\zeta^{1+n_2-n_3}+\zeta^{1-n_2+n_3}+\zeta^{n_3+2n_2}\\
&+\zeta^{-n_3-2n_2+2}+\zeta^{-n_2-2n_3+2}+\zeta^{n_2+2n_3}\\
={}&\zeta^{1+n_2-n_3}+\zeta^{2+2n_2-2n_3}+\zeta^{n_3+2n_2}+\zeta^{-n_3-2n_2+2}
 +\zeta^{-n_2-2n_3+1}+\zeta^{n_2+2n_3}\\
={}&2\cos\left( \frac{2\pi (1+n_2-n_3)}{3}\right)+2\cos\left( \frac{2\pi (2n_2+n_3)}{3}\right)\\
&+2\cos\left( \frac{2\pi (1-n_2+n_3)}{3}\right)\\
={}&0.
\end{align*}

If $n_1 \in 2+3\mathbb{Z}$.
\begin{align*}
\tilde c_{n_1-n_2-n_3,n_2,n_3}={}&\zeta^{2+n_2-n_3}+\zeta^{2-n_2+n_3}+\zeta^{n_3+2n_2}
 +\zeta^{-n_3-2n_2+1}+\zeta^{-n_2-2n_3+1}+\zeta^{n_2+2n_3}\\
={}&\zeta^{2+n_2-n_3}+\zeta^{1+2n_2-2n_3}+\zeta^{n_3+2n_2}
 +\zeta^{-n_3-2n_2+1}+\zeta^{-n_2-2n_3+2}+\zeta^{n_2+2n_3}\\
={}&2\cos\left( \frac{2\pi (2+n_2-n_3)}{3}\right)+2\cos\left( \frac{2\pi (2n_2+n_3)}{3}\right)\\
&+2\cos\left( \frac{2\pi (2-n_2+n_3)}{3}\right)\\
={}&0.\tag*{\qed}
\end{align*}
\renewcommand{\qed}{}
\end{proof}

\begin{Lemma}
Let be $u_1$, $u_2$, $u_3$ the canonical coordinates of quantum cohomology as function of the Saito flat coordinates $t^1$, $t^2$, $t^3$, i.e.,
\begin{gather}\label{canonical coordinates lemma appendix}
u_k=t^1+\frac{1}{t^3}\sum_{n=1}^{\infty} A_n^k \bigl( Q^{\frac{1}{3}}t^3 \bigr)^n.
\end{gather}
Then, the coefficients $A_{k}$ are given explicitly by
\begin{gather*}
\tilde A_{3n}=\frac{n^2N_n}{(3n-1)! },\\
\sum_{n_2=2}^{3n}3\cos\left( \frac{2\pi (n_2-1)}{3}\right) \tilde A_{3n-n_2+1} \tilde A_{n_2-1}=\bigl(6-15n-9n^2 \bigr)\frac{N_n}{(3n-1)!},\\
\sum_{n_2=1}^{3n-2}\sum_{n_3=1}^{3n-n_2-1}3\cos\left( \frac{2\pi (n_2+2n_3 )}{3} \right) \tilde A_{3n-n_2-n_3} \tilde A_{n_2} \tilde A_{n_3}\\
\qquad=\bigl(54-243n +243n^2\bigr)\frac{N_n}{(3n-1)!}+\delta_{n} ,
\end{gather*}
where
\begin{gather*}
A_n^k=\tilde A_n \bigl( {\rm e}^{\frac{2\pi {\rm i}}{3}}\bigr)^{nk},
\qquad
\delta_n=
\begin{cases}
0& \text{if } n=1,\\
 \tilde\delta_n & \text{otherwise},
\end{cases}
\\
\tilde\delta_n=\sum_{n_2=2}^{n}\frac{\bigl(6(n_2-1)-3(n-n_2+1)(n_2-1)^2 \bigr)}{(3n-3n_2+2)!(3n_2-4)!} N_{n-n_2+1}N_{n_2-1}\\
\phantom{\tilde\delta_n=}{}+\sum_{n_2=2}^{n}\frac{(-4(n-n_2+1)(n_2-1) )}{(3n-3n_2+2)!(3n_2-4)!} N_{n-n_2+1}N_{n_2-1}\\
\phantom{\tilde\delta_n=}{}+\sum_{n_2=2}^{n}\frac{\bigl(-9(n_1-n_2+1)^2(n_2-1)^2 \bigr)}{(3n-3n_2+2)!(3n_2-4)!} N_{n-n_2+1}N_{n_2-1}.
\end{gather*}
\end{Lemma}


\begin{proof}
The canonical coordinates $u_1$, $u_2$, $u_3$ can be written as
\begin{gather}\label{canonical coordinates in zi}
u_i=t^1+\frac{9+\Phi^{\prime\prime}-z_i}{t^3},
\end{gather}
where $z_i$ is are the roots of
$
 (z-z_1)(z-z_2)(z-z_3)=
 z^3-s_1z^2+s_2z-s_3=0$,
where
\begin{gather}
s_1=z_1+z_2+z_3=27+2\Phi^{\prime\prime},\nonumber\\
s_2=z_1z_2+z_2z_3+z_1z_3=243+6\Phi-15\Phi^{\prime}+27\Phi^{\prime\prime}+\bigl(\Phi^{\prime\prime}\bigr)^2,\nonumber\\
s_3=z_1z_2z_3=\bigl(27+2\Phi^{\prime}-3\Phi^{\prime\prime}\bigr)^2.\label{canonical coordinates 2 in appendix}
\end{gather}
Define
\[
f_k:=\frac{1}{t^3}\sum_{n=1}^{\infty} A_n^k \bigl( Q^{\frac{1}{3}}t^3 \bigr)^n.
\]
Substituting \eqref{canonical coordinates in zi} in \eqref{canonical coordinates lemma appendix},
\begin{gather}\label{relation zi fi}
z_i=9+\Phi^{\prime\prime}-f_i.
\end{gather}
Substituting \eqref{relation zi fi} in \eqref{canonical coordinates 2 in appendix}, we obtain
\begin{gather}
f_1+f_2+f_3=\Phi^{\prime\prime},\qquad
f_1f_2+f_2f_3+f_1f_3=6\Phi-15\Phi^{\prime}-9\Phi^{\prime\prime},\nonumber\\
f_1f_2f_3=54\Phi-243\Phi^{\prime}+243\Phi^{\prime\prime}+6\Phi\Phi^{\prime\prime}- 4{\Phi^{\prime}}^2-3\Phi^{\prime}\Phi^{\prime\prime}-9{\Phi^{\prime\prime}}^2.\label{relation fi generating function phi and derivatives}
\end{gather}
Using \eqref{generating function of Gromov Witten CP2} in the first two equation of the right-hand side of \eqref{relation fi generating function phi and derivatives}, we have
\begin{gather}
\Phi^{\prime\prime}=\sum_{n=1}^{\infty}n^2 \frac{N_n}{(3n-1)!} \bigl(Q \bigl(t^3\bigr)^3\bigr)^n ,\nonumber\\
6\Phi-15\Phi^{\prime}-9\Phi^{\prime\prime}=\sum_{n=1}^{\infty}\bigl(6-15n-9n^2 \bigr)\frac{N_n}{(3n-1)!}\bigl(Q \bigl(t^3\bigr)^3\bigr)^n.\label{relation fi generating function phi and derivatives right hand side gromov witten}
\end{gather}
The third equation of the right-hand side of \eqref{relation fi generating function phi and derivatives} is bit more involved
\begin{gather*}
54\Phi-243\Phi^{\prime}+243\Phi^{\prime\prime}+6\Phi\Phi^{\prime\prime}-3\Phi^{\prime}\Phi^{\prime\prime}-4\bigl(\Phi^{\prime}\bigr)^2-9\bigl(\Phi^{\prime\prime}\bigr)^2\nonumber\\
\qquad=\sum_{n=1}^{\infty}\bigl(54-243n+243n^2 \bigr)\frac{N_n}{(3n-1)!}\bigl(Q \bigl(t^3\bigr)^3\bigr)^n\nonumber\\
\phantom{\qquad=}{}+\sum_{n_1=1}^{\infty}\sum_{n_2=1}^{\infty}\bigl(6n_2-3n_1n_2^2-4n_1n_2-9n_1^2n_2^2 \bigr)
\frac{N_{n_1}}{(3n_1-1)!} \frac{N_{n_2}}{(3n_2-1)!}\bigl(Q \bigl(t^3\bigr)^3\bigr)^{n_1+n_2}.%\label{relation fi generating function phi and derivatives}
\end{gather*}
Using the following double infinite sum identity
\[
\sum_{k_1=1}^{\infty}\sum_{k_2=1}^{\infty} C_{k_1,k_2}=\sum_{k_1=2}^{\infty} \Bigg( \sum_{k_2=2}^{k_1} C_{k_1-k_2+1,k_2-1} \Bigg)
\]
in the equation \eqref{relation fi generating function phi and derivatives},
\begin{gather}
554\Phi-243\Phi^{\prime}+243\Phi^{\prime\prime}+6\Phi\Phi^{\prime\prime}- 4{\Phi^{\prime}}^2-3\Phi^{\prime}\Phi^{\prime\prime}-9{\Phi^{\prime\prime}}^2\nonumber\\
\qquad=\sum_{n=1}^{\infty}\bigl(54-243n +243n^2\bigr)\frac{N_n}{(3n-1)!}\bigl(Q \bigl(t^3\bigr)^3\bigr)^n\nonumber\\
\phantom{\qquad=}{}+\sum_{n_1=2}^{\infty}\left(\sum_{n_2=2}^{n_1}\frac{\bigl(6(n_2-1)-3(n_1-n_2+1)(n_2-1)^2 \bigr)}{(3n_1-3n_2+2)!(3n_2-4)!} N_{n_1-n_2+1}N_{n_2-1}\right)\bigl(Q \bigl(t^3\bigr)^3\bigr)^{n_1}\nonumber\\
\phantom{\qquad=}{}+\sum_{n_1=2}^{\infty}\left(\sum_{n_2=2}^{n_1}\frac{(-4(n_1-n_2+1)(n_2-1) )}{(3n_1-3n_2+2)!(3n_2-4)!} N_{n_1-n_2+1}N_{n_2-1}\right)\bigl(Q \bigl(t^3\bigr)^3\bigr)^{n_1}\nonumber\\
\phantom{\qquad=}{}+\sum_{n_1=2}^{\infty}\left(\sum_{n_2=2}^{n_1}\frac{\bigl(-9(n_1-n_2+1)^2(n_2-1)^2 \bigr)}{(3n_1-3n_2+2)!(3n_2-4)!} N_{n_1-n_2+1}N_{n_2-1}\right)\bigl(Q \bigl(t^3\bigr)^3\bigr)^{n_1}.\label{relation fi generating function phi and derivatives part 2}
\end{gather}
On another hand, setting
\smash{$
A_n^k=\tilde A_n \bigl( {\rm e}^{\frac{2\pi {\rm i}}{3}}\bigr)^{nk}
$}
and using the first equation of \eqref{exp 2pi 3 id}, we have
the first equation of the left-hand side of \eqref{relation fi generating function phi and derivatives} can be written as
\begin{align}
f_1+f_2+f_3&=\sum_{n=1}^{\infty} \tilde A_n \bigl( \bigl( {\rm e}^{\frac{2\pi {\rm i}}{3}}\bigr)^{n}+\bigl( {\rm e}^{\frac{2\pi {\rm i}}{3}}\bigr)^{2n}+\bigl( {\rm e}^{\frac{2\pi {\rm i}}{3}}\bigr)^{3n}\bigr)\bigl(Q^{\frac{1}{3}}t^3\bigr)^{n}\nonumber\\
&=\sum_{n=1}^{\infty} \tilde A_n \left( 1+\cos\left(\frac{2\pi n}{3}\right)\right)\bigl(Q^{\frac{1}{3}}t^3\bigr)^{n}
=\sum_{n=1}^{\infty} 3\tilde A_{3n} \bigl(Q \bigl(t^3\bigr)^3\bigr)^{n}.\label{relation fi generating function phi and derivatives left hand side}
\end{align}
Comparing equation \eqref{relation fi generating function phi and derivatives left hand side} with \eqref{relation fi generating function phi and derivatives right hand side gromov witten} and \eqref{relation fi generating function phi and derivatives}, we have
\smash{$
\tilde A_{3n}=\frac{n^2N_n}{3(3n-1)!}$}.

The second equation of the left-hand side of \eqref{relation fi generating function phi and derivatives} can be written as
\begin{align}
f_1f_2+f_2f_3+f_1f_3&=\sum_{n_1=1}^{\infty}\sum_{n_2=1}^{\infty} \tilde A_{n_1} \tilde A_{n_2} b_{n_1,n_2}\bigl(Q^{\frac{1}{3}}t^3\bigr)^{n_1+n_2}\nonumber\\
&=\sum_{n_1=2}^{\infty}\left(\sum_{n_2=2}^{n_1} \tilde A_{n_1-n_2+1} \tilde A_{n_2-1} b_{n_1-n_2+1,n_2-1}\right)\bigl(Q^{\frac{1}{3}}t^3\bigr)^{n_1},\label{relation fi generating function phi and derivatives left hand side part 2}
\end{align}
where
\smash{$
 b_{n_1,n_2}= \bigl( {\rm e}^{\frac{2\pi {\rm i}}{3}}\bigr)^{n_1+2n_2}+\bigl( {\rm e}^{\frac{2\pi {\rm i}}{3}}\bigr)^{n_1}+\bigl( {\rm e}^{\frac{2\pi {\rm i}}{3}}\bigr)^{2n_1}
$}
using Lemma \ref{lemma of exp 2pi 3} in the equation \eqref{relation fi generating function phi and derivatives left hand side part 2}
\begin{gather}
\sum_{n_1=2}^{\infty}\left(\sum_{n_2=2}^{n_1} \tilde A_{n_1-n_2+1} \tilde A_{n_2-1} b_{n_1-n_2+1,n_2-1}\right)\bigl(Q^{\frac{1}{3}}t^3\bigr)^{n_1}\nonumber\\
\qquad=\sum_{n_1=1}^{\infty}\left(\sum_{n_2=2}^{3n_1}3\cos\left( \frac{2\pi (n_2-1)}{3}\right) \tilde A_{3n_1-n_2+1} \tilde A_{n_2-1} \right)\bigl(Q^{\frac{1}{3}}t^3\bigr)^{3n_1}\nonumber\\
\qquad=\sum_{n_1=1}^{\infty}\left(\sum_{n_2=2}^{3n_1}3\cos\left( \frac{2\pi (n_2-1)}{3}\right) \tilde A_{3n_1-n_2+1} \tilde A_{n_2-1} \right)\bigl(Q\bigl(t^3\bigr)^3\bigr)^{n_1}.\label{relation fi generating function phi and derivatives left hand side part 3}
\end{gather}
Comparing \eqref{relation fi generating function phi and derivatives left hand side part 3} with \eqref{relation fi generating function phi and derivatives right hand side gromov witten},
\[%\label{2 relation final form Gromov Witten}
\sum_{n_2=2}^{3n}3\cos\left( \frac{2\pi (n_2-1)}{3}\right) \tilde A_{3n-n_2+1} \tilde A_{n_2-1}=\bigl(6-15n-9n^2 \bigr)\frac{N_n}{(3n-1)!}.
\]
	Using the triple infinite sum identity was used in
\[
\sum_{k_1=1}^{\infty}\sum_{k_2=1}^{\infty} \sum_{k_3=1}^{\infty} C_{k_1,k_2,k_3}=\sum_{k_1=3}^{\infty} \Bigg(\sum_{k_2=1}^{k_1-2} \sum_{k_3=1}^{k_1-k_2-1} C_{k_1-k_2-k_3,k_2,k_3} \Bigg).
\]
The third equation of the left-hand side of \eqref{relation fi generating function phi and derivatives} can be written as
\begin{align*}%\label{relation fi generating function phi and derivatives left hand side part 3 almost final}
f_1f_2f_3&=\sum_{n_1=1}^{\infty}\sum_{n_2=1}^{\infty}\sum_{n_3=1}^{\infty} \tilde A_{n_1} \tilde A_{n_2} \tilde A_{n_3} c_{n_1,n_2,n_3}\bigl(Q^{\frac{1}{3}}t^3\bigr)^{n_1+n_2+n_3}\\
&=\sum_{n_1=3}^{\infty}\left(\sum_{n_2=1}^{n_1-2}\sum_{n_3=1}^{n_1-n_2-1} \tilde A_{n_1-n_2-n_3} \tilde A_{n_2} \tilde A_{n_3} c_{n_1-n_2-n_3,n_2,n_3}\right)\bigl(Q^{\frac{1}{3}}t^3\bigr)^{n_1},
\end{align*}
where
\smash{$
c_{n_1,n_2,n_3}=\bigl( {\rm e}^{\frac{2\pi {\rm i}}{3}} \bigr)^{n_1+2n_2+3n_3}$}.
Then
\begin{gather}
\sum_{n_1=3}^{\infty}\left(\sum_{n_2=1}^{n_1-2}\sum_{n_3=1}^{n_1-n_2-1} \tilde A_{n_1-n_2-n_3} \tilde A_{n_2} \tilde A_{n_3} c_{n_1-n_2-n_3,n_2,n_3}\right)\bigl(Q^{\frac{1}{3}}t^3\bigr)^{n_1}\nonumber\\
\qquad=\sum_{n_1=3}^{\infty}\left(\sum_{n_2=1}^{n_1-2}\sum_{n_3=1}^{n_1-n_2-1} \tilde A_{n_1-n_2-n_3} \tilde A_{n_2} \tilde A_{n_3} \bigl( {\rm e}^{\frac{2\pi {\rm i}}{3}} \bigr)^{n_1+n_2+2n_3}\right)\bigl(Q^{\frac{1}{3}}t^3\bigr)^{n_1}\nonumber\\
\qquad=\sum_{n_1=1}^{\infty}\left(\sum_{n_2=1}^{3n_1-2}\sum_{n_3=1}^{3n_1-n_2-1}3\cos\left( \frac{2\pi (n_2+2n_3 )}{3} \right) \tilde A_{3n_1-n_2-n_3} \tilde A_{n_2} \tilde A_{n_3} \right)\bigl(Q^{\frac{1}{3}}t^3\bigr)^{3n_1}\nonumber\\
\qquad=\sum_{n_1=1}^{\infty}\left(\sum_{n_2=1}^{3n_1-2}\sum_{n_3=1}^{3n_1-n_2-1}3\cos\left( \frac{2\pi (n_2+2n_3 )}{3} \right) \tilde A_{3n_1-n_2-n_3} \tilde A_{n_2} \tilde A_{n_3} \right)\bigl(Q\bigl(t^3\bigr)^3\bigr)^{n_1}.\label{relation fi generating function phi and derivatives left hand side part 3 almost final part 1}
\end{gather}
Comparing \eqref{relation fi generating function phi and derivatives left hand side part 3 almost final part 1} with \eqref{relation fi generating function phi and derivatives part 2},
\begin{gather*}%\label{relation fi generating function phi and derivatives left hand side part 3 almost final part 1 final part ever}
\sum_{n_2=1}^{3n-2}\sum_{n_3=1}^{3n-n_2-1}3\cos\left( \frac{2\pi (n_2+2n_3 )}{3} \right) \tilde A_{3n-n_2-n_3} \tilde A_{n_2} \tilde A_{n_3}\\
\qquad=\bigl(54-243n+243n^2 \bigr)\frac{N_n}{(3n-1)!}+\delta_{n} ,
\end{gather*}
where
\begin{gather*}
\delta_n=
\begin{cases}
0& \text{if } n=1,\\
 \tilde\delta_n & \text{otherwise},
\end{cases}
\\
\tilde\delta_n=\sum_{n_2=2}^{n}\frac{\bigl(6(n_2-1)-3(n-n_2+1)(n_2-1)^2 \bigr)}{(3n-3n_2+2)!(3n_2-4)!} N_{n-n_2+1}N_{n_2-1}\\
\phantom{\tilde\delta_n=}{}+\sum_{n_2=2}^{n}\frac{(-4(n-n_2+1)(n_2-1) )}{(3n-3n_2+2)!(3n_2-4)!} N_{n-n_2+1}N_{n_2-1}\\
\phantom{\tilde\delta_n=}{}+\sum_{n_2=2}^{n}\frac{(-9(n_1-n_2+1)^2(n_2-1)^2 )}{(3n-3n_2+2)!(3n_2-4)!} N_{n-n_2+1}N_{n_2-1}.\tag*{\qed}
\end{gather*} \renewcommand{\qed}{}
\end{proof}

\begin{Corollary}
Let be $u_1$, $u_2$, $u_3$ the canonical coordinates of quantum cohomology as function of the Saito flat coordinates $t^1$, $t^2$, $t^3$, i.e.,
\begin{gather*}%\label{canonical coordinates lemma appendix}
u_k=t^1+\frac{1}{t^3}\sum_{n=1}^{\infty} A_n^k \bigl( Q^{\frac{1}{3}}t^3 \bigr)^n.
\end{gather*}
Then, the coefficients $A_{k}$ are given recursively by
\begin{gather*}
\tilde A_{3n}=\frac{n^2N_n}{(3n-1)! },\\
 \tilde A_{3n-1}=\frac{1}{3 \tilde A_{1}}\left[\bigl(6-15n-9n^2 \bigr)\frac{N_n}{(3n-1)!}-\sum_{n_2=3}^{3n-1}3\cos\left( \frac{2\pi (n_2-1)}{3}\right) \tilde A_{3n-n_2+1} \tilde A_{n_2-1}\right],\\
 \tilde A_{3n-2}=\frac{1}{9\tilde A_1^2}\left[\bigl(54-243n +243n^2\bigr)\frac{N_n}{(3n-1)!}+\delta_{n} \right.\\
 \phantom{ \tilde A_{3n-2}=}{}\left.-\sum_{n_3=2}^{3n-2}3\cos\left( \frac{2\pi (1+2n_3 )}{3} \right) \tilde A_{3n-1-n_3} \tilde A_{1} \tilde A_{n_3}\right]\\
\phantom{ \tilde A_{3n-2}=}{}-\frac{1}{9\tilde A_1^2}\left[\sum_{n_2=2}^{3n-3}\sum_{n_3=1}^{3n-n_2-1}3\cos\left( \frac{2\pi (n_2+2n_3 )}{3} \right) \tilde A_{3n-n_2-n_3} \tilde A_{n_2} \tilde A_{n_3}\right],
\end{gather*}
where
\begin{gather*}
A_n^k=\tilde A_n \bigl( {\rm e}^{\frac{2\pi {\rm i}}{3}}\bigr)^{nk}
,\qquad
\delta_n=
\begin{cases}
0& \text{if } n=1,\\
 \tilde\delta_n & \text{otherwise},
\end{cases}
\\
\tilde\delta_n=\sum_{n_2=2}^{n}\frac{\bigl(6(n_2-1)-3(n-n_2+1)(n_2-1)^2 \bigr)}{(3n-3n_2+2)!(3n_2-4)!} N_{n-n_2+1}N_{n_2-1}\\
\phantom{\tilde\delta_n=}{}+\sum_{n_2=2}^{n}\frac{(-4(n-n_2+1)(n_2-1) )}{(3n-3n_2+2)!(3n_2-4)!} N_{n-n_2+1}N_{n_2-1}\\
\phantom{\tilde\delta_n=}{}+\sum_{n_2=2}^{n}\frac{\bigl(-9(n_1-n_2+1)^2(n_2-1)^2 \bigr)}{(3n-3n_2+2)!(3n_2-4)!} N_{n-n_2+1}N_{n_2-1}.
\end{gather*}
\end{Corollary}

\subsection{Coefficients of the cross ratio function}\label{Coefficients of the cross ratio function}

\begin{Lemma}
Let the cross ratio function be defined by the formula
\begin{gather}
f\bigl(Q^{\frac{1}{3}}t^3\bigr)=\frac{u_3-u_1}{u_2-u_1} =\sum_{n=0}^{\infty} f_n \bigl(Q^{\frac{1}{3}}t^3\bigr)^n.\label{identity cross ratio qt3}
\end{gather}
Then, the Taylor series of \eqref{identity cross ratio qt3} is given by
\[%\label{identity cross ratio qt3 2}
f\bigl(Q^{\frac{1}{3}}t^3\bigr)=\sum_{n=0}^{\infty} \sum_{m=0}^{n}\left(\frac{A_{n+1-m}^3-A_{n+1-m}^1}{A_1^2-A_1^1}\right)C_{m,-1}(y_1,y_2,\dots,y_m) \bigl( Q^{\frac{1}{3}}t^3 \bigr)^{n},
\]
where
\smash{$
y_n=\frac{A_{n+1}^2-A_{n+1}^1}{A_1^2-A_1^1}$},
and $A_n$ is defined in \eqref{canonical coordinates lemma1}.
\end{Lemma}

\begin{proof}
The coefficients $f_n$ \eqref{identity cross ratio qt3} is obtained by substituting \eqref{canonical coordinates lemma1} in \eqref{identity cross ratio qt3} and using Theorem~\ref{theorem potential partition polynomials}. Indeed,
\begin{align*}%\label{identity cross ratio qt3 coed}
\frac{u_3-u_1}{u_2-u_1}&= \frac{\sum_{n=1}^{\infty} \bigl(A_n^3-A_n^1\bigr) \bigl( Q^{\frac{1}{3}}t^3 \bigr)^n}{\sum_{n=1}^{\infty} \bigl(A_n^2-A_n^1\bigr) \bigl( Q^{\frac{1}{3}}t^3 \bigr)^n}\\
&=\left(\sum_{n=0}^{\infty} \left(\frac{A_{n+1}^3-A_{n+1}^1}{A_1^2-A_1^1}\right) \bigl( Q^{\frac{1}{3}}t^3 \bigr)^n\right)\left(1+\sum_{n=1}^{\infty} \left(\frac{A_{n+1}^2-A_{n+1}^1}{A_1^2-A_1^1}\right) \bigl( Q^{\frac{1}{3}}t^3 \bigr)^n\right)^{-1}\\
&=\left(\sum_{n=0}^{\infty} \left(\frac{A_{n+1}^3-A_{n+1}^1}{A_1^2-A_1^1}\right) \bigl( Q^{\frac{1}{3}}t^3 \bigr)^n\right)\left(\sum_{n=0}^{\infty} C_{n,-1}(y_1,y_2,\dots,y_n) \bigl( Q^{\frac{1}{3}}t^3 \bigr)^n\right)\\
&=\sum_{n=0}^{\infty} \sum_{m=0}^{\infty}\left(\frac{A_{n+1}^3-A_{n+1}^1}{A_1^2-A_1^1}\right)C_{m,-1}(y_1,y_2,\dots,y_m) \bigl( Q^{\frac{1}{3}}t^3 \bigr)^{n+m}\\
&=\sum_{n=0}^{\infty} \sum_{m=0}^{n}\left(\frac{A_{n+1-m}^3-A_{n+1-m}^1}{A_1^2-A_1^1}\right)C_{m,-1}(y_1,y_2,\dots,y_m) \bigl( Q^{\frac{1}{3}}t^3 \bigr)^{n},
\end{align*}
where
\smash{$
y_n=\frac{A_{n+1}^2-A_{n+1}^1}{A_1^2-A_1^1}$},
and the following Double sum identity was used:
\[
\sum_{n=0}^{\infty}\sum_{m=0}^{\infty} C_{n,m}=\sum_{n=0}^{\infty}\sum_{m=0}^{n} C_{n-m,m}.
\]
Lemma proved.
\end{proof}

\subsection{Coefficients of modular lambda function}\label{Coefficients of Modular lambda function}

Recall the Eisenstein series $E_2(\tau)$, $E_4(\tau)$, $E_6(\tau)$ at the point $\tau={\rm e}^{\frac{2\pi {\rm i}}{3}}$ have the following values:
\begin{gather}\label{special values of Eisenstein series in the equiharmonic lattice }
E_2\bigl({\rm e}^{\frac{2\pi {\rm i}}{3}}\bigr)=\frac{2\sqrt{3}}{\pi}, \qquad E_4\bigl({\rm e}^{\frac{2\pi {\rm i}}{3}}\bigr)=0, \qquad E_6\bigl({\rm e}^{\frac{2\pi {\rm i}}{3}}\bigr)=\frac{\bigl(\Gamma\bigl(\frac{1}{3}\bigr)\bigr)^{18}}{4^6}.
\end{gather}
Moreover, there exist the following relation between $x$, $x^{\prime}$ and $\Delta$:
\begin{gather}
\frac{2^8}{3^3}(2\pi )^{-12}\Delta(\tau)=\frac{(x^{\prime})^6}{x^4(x-1)^4}.\label{main relationship between x,xprime and delta}
\end{gather}
The special values of $x$ at $\tau={\rm e}^{\frac{2\pi {\rm i}}{3}}$ is
\begin{gather}
x\bigl({\rm e}^{\frac{2\pi {\rm i}}{3}}\bigr)={\rm e}^{\frac{\pi {\rm i}}{3}}.\label{special value of x in the equiharmonic lattice}
\end{gather}
Substituting \eqref{special value of x in the equiharmonic lattice}, \eqref{special values of Eisenstein series in the equiharmonic lattice } in \eqref{main relationship between x,xprime and delta}, we compute $x^{\prime}\bigl({\rm e}^{\frac{2\pi {\rm i}}{3}}\bigr)$. Computing
\begin{gather}
\frac{2^8}{3^3}(2\pi )^{-12}\frac{{\rm d}^n\log\Delta(\tau)}{{\rm d}\tau^n}=\frac{{\rm d}^n}{{\rm d}\tau^n}\left(\frac{(x^{\prime})^6}{x^4(x-1)^4}\right)n,\label{main relationship between x,xprime and delta derivatives }
\end{gather}
we obtain in the left-hand side \eqref{main relationship between x,xprime and delta derivatives } a polynomial expression in terms of $E_2$, $E_4$, $E_6$ and in the right-hand side a rational expression of \smash{$\frac{{\rm d}^nx(\tau)}{{\rm d}\tau^n}$}. Then, we can derive \smash{$\frac{{\rm d}^nx(\tau)}{{\rm d}\tau^n}$} for any order $n$ by using the recursive relation \eqref{main relationship between x,xprime and delta derivatives }.
